% Options for packages loaded elsewhere
\PassOptionsToPackage{unicode}{hyperref}
\PassOptionsToPackage{hyphens}{url}
\PassOptionsToPackage{dvipsnames,svgnames*,x11names*}{xcolor}
%
\documentclass[
  11pt,
]{article}
\usepackage{lmodern}
\usepackage{amssymb,amsmath}
\usepackage{ifxetex,ifluatex}
\ifnum 0\ifxetex 1\fi\ifluatex 1\fi=0 % if pdftex
  \usepackage[T1]{fontenc}
  \usepackage[utf8]{inputenc}
  \usepackage{textcomp} % provide euro and other symbols
\else % if luatex or xetex
  \usepackage{unicode-math}
  \defaultfontfeatures{Scale=MatchLowercase}
  \defaultfontfeatures[\rmfamily]{Ligatures=TeX,Scale=1}
  \setmainfont[]{DejaVu Serif}
  \setsansfont[]{DejaVu Sans}
  \setmonofont[]{DejaVu Sans Mono}
\fi
% Use upquote if available, for straight quotes in verbatim environments
\IfFileExists{upquote.sty}{\usepackage{upquote}}{}
\IfFileExists{microtype.sty}{% use microtype if available
  \usepackage[]{microtype}
  \UseMicrotypeSet[protrusion]{basicmath} % disable protrusion for tt fonts
}{}
\makeatletter
\@ifundefined{KOMAClassName}{% if non-KOMA class
  \IfFileExists{parskip.sty}{%
    \usepackage{parskip}
  }{% else
    \setlength{\parindent}{0pt}
    \setlength{\parskip}{6pt plus 2pt minus 1pt}}
}{% if KOMA class
  \KOMAoptions{parskip=half}}
\makeatother
\usepackage{xcolor}
\IfFileExists{xurl.sty}{\usepackage{xurl}}{} % add URL line breaks if available
\IfFileExists{bookmark.sty}{\usepackage{bookmark}}{\usepackage{hyperref}}
\hypersetup{
  pdftitle={Methods: High-Speed Imaging Analysis of Plasma Arc Dynamics},
  colorlinks=true,
  linkcolor=RoyalBlue,
  filecolor=Maroon,
  citecolor=Blue,
  urlcolor=RoyalBlue,
  pdfcreator={LaTeX via pandoc}}
\urlstyle{same} % disable monospaced font for URLs
\usepackage[margin=2.5cm]{geometry}
\usepackage{longtable,booktabs}
% Correct order of tables after \paragraph or \subparagraph
\usepackage{etoolbox}
\makeatletter
\patchcmd\longtable{\par}{\if@noskipsec\mbox{}\fi\par}{}{}
\makeatother
% Allow footnotes in longtable head/foot
\IfFileExists{footnotehyper.sty}{\usepackage{footnotehyper}}{\usepackage{footnote}}
\makesavenoteenv{longtable}
\setlength{\emergencystretch}{3em} % prevent overfull lines
\providecommand{\tightlist}{%
  \setlength{\itemsep}{0pt}\setlength{\parskip}{0pt}}
\setcounter{secnumdepth}{-\maxdimen} % remove section numbering
\usepackage{microtype}
\setlength{\emergencystretch}{3em}
\usepackage{booktabs}
\usepackage{longtable}

\title{Methods: High-Speed Imaging Analysis of Plasma Arc Dynamics}
\usepackage{etoolbox}
\makeatletter
\providecommand{\subtitle}[1]{% add subtitle to \maketitle
  \apptocmd{\@title}{\par {\large #1 \par}}{}{}
}
\makeatother
\subtitle{Unified analysis for both electrode geometries (CO₂
J-electrode and N₂ blade)}
\author{}
\date{21 June 2026}

\begin{document}
\maketitle

{
\hypersetup{linkcolor=}
\setcounter{tocdepth}{2}
\tableofcontents
}
\begin{quote}
\textbf{Purpose.} A single, paper-oriented description of the
gliding-arc image analysis used for \textbf{both} reactor geometries,
with the shared core stated once and the geometry-specific pieces placed
side by side. Use this to talk through the methods end to end or to
draft a manuscript Methods section for either experiment.

Last revised: 2026-06-21. Detailed per-experiment companions:
\texttt{METHODS.md} (CO₂ J-electrode) and \texttt{METHODS\_BLADE.md} (N₂
blade). Where this document and a companion disagree, the companion is
authoritative for its experiment.
\end{quote}

\begin{center}\rule{0.5\linewidth}{0.5pt}\end{center}

\hypertarget{1-scope-two-experiments-one-analysis-stack}{%
\subsection{1. Scope: two experiments, one analysis
stack}\label{1-scope-two-experiments-one-analysis-stack}}

Two high-speed discharge experiments are analysed with a \textbf{common}
image- analysis pipeline. They differ only in gas, electrode geometry,
calibration, and operating sweep; every quantitative measurement
(detection, skeletonisation, event segmentation, velocity, modes,
uncertainty) is computed by the \textbf{same} code.

\begin{longtable}[]{@{}lll@{}}
\toprule
& \textbf{CO₂ J-electrode} & \textbf{N₂ blade}\tabularnewline
\midrule
\endhead
Gas & CO₂ & N₂\tabularnewline
Electrodes & hooked \textbf{finite} J-segments & straight
\textbf{infinite} blade lines (25° half-angle wedge)\tabularnewline
Operating points & sweep: nozzle 0.4/0.8/1.2 mm × flow 1000/1500
mL·min⁻¹, with a repeat & single: 0.6 mm, 1.5 SLPM, two reactor
orientations\tabularnewline
Frame rate & per-recording, auto-detected (≈ 200 kHz) & 200 kHz (5
µs/frame)\tabularnewline
Calibration & per-recording px/mm from ≈ 2.4 mm apex gap & 8.8 px/mm (44
px = 5 mm)\tabularnewline
Orchestrator & \texttt{run\_j\_co2.py} &
\texttt{run\_11may26.py}\tabularnewline
Companion doc & \texttt{METHODS.md} &
\texttt{METHODS\_BLADE.md}\tabularnewline
\bottomrule
\end{longtable}

\textbf{Shared modules} (identical code for both):
\texttt{analyze\_blade.py} (detection + per-frame measurement),
\texttt{event\_velocity.py} (events, velocity, modes, uncertainty),
\texttt{arc\_shape.py} (normal/tangential decomposition),
\texttt{uncertainty.py} (systematic budget),
\texttt{build\_velocity\_map.py}, \texttt{fourier\_analysis.py}. The two
orchestrators differ only in the geometry/calibration they pass and, for
the J path, the finite-segment attachment gate (§5).

\begin{center}\rule{0.5\linewidth}{0.5pt}\end{center}

\hypertarget{2-physical-picture-common}{%
\subsection{2. Physical picture
(common)}\label{2-physical-picture-common}}

In both reactors two electrodes form a narrow gap (the \textbf{apex})
that widens downstream. The arc breaks down at the apex, then the gas
jet and electric field drive it to \textbf{extend/glide} along the
diverging electrodes, forming an arch whose feet stay rooted on the
electrodes. The arc then extinguishes and \textbf{re-nucleates},
producing a periodic discharge cycle. Gliding-arc theory distinguishes
two re-nucleation modes, both classified by the pipeline:

\begin{itemize}
\tightlist
\item
  \textbf{Mode A} --- full extinction, fresh re-ignition at the apex
  (minimum gap).
\item
  \textbf{Mode B} --- back-breakdown / reconnection at an intermediate
  position.
\end{itemize}

(Czernichowski 1994; Zhang et al. 2020; Fridman et al. 2005.)

\begin{center}\rule{0.5\linewidth}{0.5pt}\end{center}

\hypertarget{3-pipeline-architecture-common}{%
\subsection{3. Pipeline architecture
(common)}\label{3-pipeline-architecture-common}}

\begin{verbatim}
orchestrator (run_j_co2.py | run_11may26.py)
  Stage 1  convert    .cine → 8-bit TIFF chunks
                      • percentile-clip normalisation from sampled frames
                      • per-chunk bad-pixel correction (mandatory; see §4)
  Stage 2  analyze    analyze_blade.py per chunk → per-frame measurements.csv
  Stage 3  merge      concatenate chunk CSVs with global frame numbering
  Stage 4  events     event_velocity.py → events, velocities, Mode A/B, σ
  Stage 5  fourier    fourier_analysis.py → cycle spectra (peaks ± σ_f)   [J path]
  Stage 6  visualize  build_velocity_map.py + visualize_blade.py → figures
\end{verbatim}

QA (\texttt{qa\_blade.py}) is run on representative chunks before a full
run to confirm detection quality. All resolved parameters and provenance
are written to \texttt{run\_metadata.json} (§11).

\begin{center}\rule{0.5\linewidth}{0.5pt}\end{center}

\hypertarget{4-image-conversion--preprocessing-common}{%
\subsection{4. Image conversion \& preprocessing
(common)}\label{4-image-conversion--preprocessing-common}}

\texttt{.cine} frames (12-bit, 0--4095) are converted to 8-bit TIFF
chunks with percentile-clip normalisation estimated from sampled frames.
\textbf{Bad-pixel correction is mandatory}: the sensor carries a fixed
grid of \textasciitilde128 stuck/hot pixels that saturate every frame;
\texttt{correct\_bad\_pixels\_stack} removes them during conversion. Any
path reading raw frames without this step registers false detections on
the hot-pixel grid (\texttt{qa\_blade.py} reports
\texttt{persistent\_saturated\_px}, which is 0 on correctly converted
chunks).

\textbf{Intensity convention.} Quantitative brightness is measured on
\textbf{raw DN} (0--255); cross-frame comparisons use
\texttt{integrated\_intensity}, never a per-frame normalised value.

\begin{center}\rule{0.5\linewidth}{0.5pt}\end{center}

\hypertarget{5-electrode-geometry--spatial-calibration}{%
\subsection{5. Electrode geometry \& spatial
calibration}\label{5-electrode-geometry--spatial-calibration}}

This is the \textbf{principal divergence} between the two experiments.
The shared requirement is that a physically meaningful bridging arc
roots on \textbf{both} electrodes; the geometry determines how "rooted
on the electrode" is tested and how px/mm is set.

\hypertarget{shared-pxmm-from-a-known-gap-and-the-apex-anchor}{%
\subsubsection{Shared: px/mm from a known gap, and the apex
anchor}\label{shared-pxmm-from-a-known-gap-and-the-apex-anchor}}

Spatial scale comes from a known physical gap measured in pixels. The
apex gap also anchors the electrode-gap \textbf{flash-cutoff row} and
the Mode-A/B reference position. Velocity uncertainty includes a
calibration term derived from this gap (§10).

\hypertarget{5a-j-electrode-coux2082-finite-hooked-segments}{%
\subsubsection{5a. J-electrode (CO₂): finite hooked
segments}\label{5a-j-electrode-coux2082-finite-hooked-segments}}

Each electrode is a \textbf{finite} curved J-segment. Two
representations are supported: straight finite blades
(\texttt{blade\_lines.npy} from the interactive
\texttt{calibrate\_j\_electrodes.py}) or curved edges
(\texttt{electrode\_edges.npz} from \texttt{detect\_j\_hooks.py}).
Attachment is enforced by an explicit \textbf{finite-segment gate}
(\texttt{-\/-finite-segments}): each arc foot must land on the real
electrode segment near the apex, not on an imaginary infinite extension.
This gate is what keeps incomplete arcs out of the event stream.

px/mm is \textbf{per-recording}, written to \texttt{j\_geometry.json},
resolved as
\texttt{-\/-px-per-mm\ (CLI)\ \textgreater{}\ j\_geometry.json\ \textgreater{}\ tiff/calibration.txt\ \textgreater{}\ 9.167\ (fallback)},
from the ≈ 2.4 mm apex gap. \emph{Known provenance issue:} the 1.2 mm
recording's \texttt{j\_geometry.json} stores
\texttt{px\_per\_mm\ =\ 8.75} while the stored gap implies 11.17; runs
use a \texttt{-\/-px-per-mm\ 11.2} override (numbers correct, stored
field stale).

\hypertarget{5b-blade-nux2082-straight-infinite-lines}{%
\subsubsection{5b. Blade (N₂): straight infinite
lines}\label{5b-blade-nux2082-straight-infinite-lines}}

Each electrode is a straight blade modelled as an \textbf{infinite
line}. Attachment is tested by the perpendicular distance of each
skeleton foot to a stored blade line (\texttt{blade\_lines.npy},
tolerance \texttt{ELEC\_TOL\_PX}). The lines are fit automatically by
\texttt{detect\_blade\_lines.py} (total-least-squares on the loci of arc
feet across clean-arch frames) or overridden interactively by
\texttt{calibrate\_electrodes.py}. px/mm is the single dataset constant
\textbf{8.8 px/mm} (44 px = 5 mm). The "25deg" filename label is the
per-blade \textbf{half-angle}; the auto-detected full opening angle is ≈
38--47°.

\hypertarget{orientation-handling-common}{%
\subsubsection{Orientation handling
(common)}\label{orientation-handling-common}}

\texttt{arc\_tip\_row} is stored as \textbf{distance from the
electrode}, increasing as the arc extends. For inverted (upward)
recordings the raw row is flipped (\texttt{H-1-row}) so downstream code
sees one consistent advancing-front axis regardless of reactor
orientation.

\begin{center}\rule{0.5\linewidth}{0.5pt}\end{center}

\hypertarget{6-arc-detection-common-fixed-dn-threshold}{%
\subsection{6. Arc detection (common: fixed DN
threshold)}\label{6-arc-detection-common-fixed-dn-threshold}}

Detection uses a \textbf{fixed DN threshold}, not adaptive Otsu --- the
arc is a small bright object on a near-black field, and per-frame
adaptive thresholds amplify noise on dark inter-stroke frames. Steps
(\texttt{analyze\_blade.binarize}): dark guard (skip frames with max
\textless{} \textasciitilde35 DN) → threshold at \texttt{bright\_thresh}
→ morphological closing (disk r = 3) → keep the single largest connected
component within \texttt{{[}min\_arc\_px,\ max\_arc\_px{]}}.

\textbf{Threshold provenance differs by experiment:}

\begin{itemize}
\tightlist
\item
  \textbf{J (CO₂):} resolved per-recording,
  \texttt{CLI\ \textgreater{}\ j\_geometry.json\ \textgreater{}\ tiff/threshold\_suggested.txt\ \textgreater{}\ 210\ (fallback)}.
  The 210 fallback (and per-recording values, e.g. 160 / 220 / 235) come
  from arc-vs-plume local-intensity QA
  (\texttt{qa\_arc\_local\_intensity\_j.py}): arc core saturates near
  255, plume below \textasciitilde150. \texttt{max\_arc\_px} is raised
  above the blade default because the wide high-flow CO₂ arc can exceed
  it.
\item
  \textbf{Blade (N₂):} a single dataset-wide constant \textbf{100 DN},
  from \texttt{threshold\_qa\_11may26.py}, since it is one operating
  condition.
\end{itemize}

Flash/saturation rejection (\texttt{arc\_filter.py}) applies the unified
validity criteria: a gap-row flash cutoff, an \texttt{arc\_width\_px}
leg, and an effective-width leg
(\texttt{eff\_width\ =\ area\ /\ length}, which fires on disconnected
double-strikes). \texttt{filter\_audit\_montage.py} produces a
reproducible accepted/rejected montage to confirm the criteria behave.

\begin{center}\rule{0.5\linewidth}{0.5pt}\end{center}

\hypertarget{7-per-frame-measurements-common-analyze_blademeasure_frame}{%
\subsection{\texorpdfstring{7. Per-frame measurements (common,
\texttt{analyze\_blade.measure\_frame})}{7. Per-frame measurements (common, analyze\_blade.measure\_frame)}}\label{7-per-frame-measurements-common-analyze_blademeasure_frame}}

For each detected frame: arc-mask \textbf{area}; \textbf{skeleton}
(Zhang--Suen) and ordered centreline (depth-first walk on the
8-connected graph; endpoints = degree-1 pixels); \textbf{arc length}
(skeleton px / px-per-mm); \textbf{FWHM width} sampled perpendicular to
the local tangent (channel-diameter proxy \textbf{plus} motion blur at
the few-µs exposure); \textbf{brightness} (mean / peak / integrated, raw
DN); \textbf{centroid}; \textbf{bounding box}; \textbf{curvature} (mean
\textbar Δθ\textbar{} along the skeleton); the two skeleton
\textbf{endpoints} (arc feet); and the \textbf{arc-front tip} (extreme
pixel in the propagation direction).

\begin{center}\rule{0.5\linewidth}{0.5pt}\end{center}

\hypertarget{8-events-velocity--modes-common-event_velocitypy}{%
\subsection{\texorpdfstring{8. Events, velocity \& modes (common,
\texttt{event\_velocity.py})}{8. Events, velocity \& modes (common, event\_velocity.py)}}\label{8-events-velocity--modes-common-event_velocitypy}}

Operates on the merged per-frame CSV; the only per-experiment inputs are
\texttt{-\/-px-per-mm} and \texttt{-\/-fps}.

\begin{quote}
\textbf{fps invariant.} fps is per-recording (auto-detected) and must
thread into \texttt{event\_velocity.py}, \texttt{qa\_restrike.py}, and
\texttt{fourier\_analysis.py}. A historic regression mis-scaled velocity
and frequency by \texttt{true\_fps\ /\ 200000}; this is now guarded by a
dimensional self-check (\texttt{test\_j\_geometry.py}: synthetic 200 kHz
vs 100 kHz → ratio exactly 2.0) covering both geometries.
\end{quote}

\textbf{Event detection.} A new primary event begins when the tip
returns to within 15 \% of the effective arc range above the electrode
row, or after a dropout

\begin{quote}
500 frames. The electrode row is the 5th-percentile
\texttt{arc\_tip\_row} of complete-arc frames;
\texttt{L\_eff\_px\ =\ p90(tip)\ −\ electrode\_row}. Retreats
\textgreater{} 50 \% of the current extension are flagged as sub-events.
\end{quote}

\textbf{Velocity.} The slope of an OLS fit to the
\textbf{advancing-front} (monotonic- envelope) tip trajectory:

\begin{verbatim}
v = slope [px/frame] × FPS / px_per_mm / 1000   (m/s)
\end{verbatim}

R² reports linearity. A centered sliding-window OLS over \textbf{N = 10}
advancing- front frames (±5; \texttt{-\/-local-vel-n}) gives a
time-resolved local velocity (equivalent to a Savitzky--Golay first
derivative). \textbf{N = 10 is canonical}: N = 5 spans only
\textasciitilde25 µs and admits short-baseline spikes; the J sensitivity
sweep confirms the headline metric is stable for N ≥ 7 (§11).

\textbf{Mode A / B.} From the dimensionless normalised nucleation
position \texttt{(tip\_row\_start\ −\ electrode\_row)\ /\ L\_eff\_px},
threshold 0.15 (self-consistent with the event detector). The
authoritative record is \texttt{re\_strikes\_qa.csv}.

\textbf{Output.} \texttt{event\_summary.csv} carries both
\texttt{vel\_ms} (whole-event average rate) and time-resolved
\texttt{local\_vel\_\{median,mean,p10,p90,peak,max,std\}}; the full
per-frame profile is in \texttt{events/event\_NNNN\_dynamics.csv}. The
two diverge for multi-phase events --- the signal that a single average
is insufficient.

\begin{center}\rule{0.5\linewidth}{0.5pt}\end{center}

\hypertarget{9-arc-shape-tracking--normal-velocity-common-arc_shapepy}{%
\subsection{\texorpdfstring{9. Arc-shape tracking \& normal velocity
(common,
\texttt{arc\_shape.py})}{9. Arc-shape tracking \& normal velocity (common, arc\_shape.py)}}\label{9-arc-shape-tracking--normal-velocity-common-arc_shapepy}}

The arc is a deformable filament: the tip regression captures only the
leading edge. \texttt{arc\_shape.py} resamples the skeleton to N
arc-length-equally-spaced points (point 0 = electrode-proximal) and, for
each consecutive frame pair within an event, decomposes each point's
displacement into \textbf{tangential} (along-arc) and \textbf{normal}
(lateral sweep) components relative to the local arc direction. Only
complete-arc frames are used. Output \texttt{arc\_shape\_velocities.csv}
with \texttt{norm\_pos\ =\ seg\_idx\ /\ (N\ −\ 1)}. This separates
gliding motion from jet-driven lateral sweep for either geometry.

\begin{center}\rule{0.5\linewidth}{0.5pt}\end{center}

\hypertarget{10-uncertainty-budget-common-uncertaintypy}{%
\subsection{\texorpdfstring{10. Uncertainty budget (common,
\texttt{uncertainty.py})}{10. Uncertainty budget (common, uncertainty.py)}}\label{10-uncertainty-budget-common-uncertaintypy}}

A single source of truth for the systematic budget the velocity stack
previously lacked: the OLS slope SE is only the random fitting error;
reported velocities must also carry the \textbf{px/mm} and \textbf{fps}
systematics. For \texttt{v\ =\ slope\ ·\ FPS\ /\ px\_per\_mm\ /\ 1000}:

\begin{verbatim}
(σ_v / |v|)²  =  (σ_v_stat / |v|)²  +  (σ_fps / FPS)²  +  (σ_pxmm / px_per_mm)²
\end{verbatim}

\begin{longtable}[]{@{}lll@{}}
\toprule
Quantity & Default & Justification\tabularnewline
\midrule
\endhead
\texttt{SIGMA\_FPS\_REL} & 0.1 \% & Phantom HS crystal ≈ 50--100 ppm;
conservative for trig-mode\tabularnewline
\texttt{SIGMA\_PXMM\_REL} & 5 \% & √2 × 1 px click / ≈ 26.56 px apex gap
≈ 5.3 \%\tabularnewline
\bottomrule
\end{longtable}

Both are relative 1-σ, overridable per-recording
(\texttt{-\/-pxmm-rel-err}, \texttt{-\/-fps-rel-err}) and logged in
\texttt{run\_metadata.json}.
\texttt{pxmm\_uncertainty\_from\_calibration()} can derive σ\_pxmm from
the measured gap. \texttt{event\_velocity.py} writes
\texttt{vel\_ms\_err\_stat}, \texttt{vel\_ms\_err\_sys},
\texttt{vel\_ms\_err\_total} (\textbf{report the total}); the legacy
\texttt{vel\_ms\_err} is retained = stat only. Frequencies obey
\texttt{σ\_f\ /\ f\ =\ σ\_fps\ /\ FPS}; Fourier peaks are annotated
\texttt{f\ ±\ σ\_f}. \texttt{recompute\_uncertainty.py} adds these
columns to existing results without re-running --- point estimates are
unchanged; only the error bars grow (≈ 6× on the J target, where the
calibration term had been omitted entirely).

This budget applies identically to both geometries.

\begin{center}\rule{0.5\linewidth}{0.5pt}\end{center}

\hypertarget{11-aggregate--robustness-analyses}{%
\subsection{11. Aggregate \& robustness
analyses}\label{11-aggregate--robustness-analyses}}

These operate across or within recordings; their applicability differs
because only the CO₂ experiment is a multi-condition sweep.

\textbf{Sensitivity sweep} (\texttt{sensitivity\_sweep.py}, Tier A;
post-hoc on \texttt{measurements\_all\_events.csv}). A parameter is
UNSTABLE if shifting it moves the median by more than
\texttt{σ\_floor\ =\ √(σ\_pxmm²\ +\ σ\_fps²)\ ×\ median\_v}. On the J
target: \texttt{local\_vel\_n} stable for N ≥ 7 (N = 5 UNSTABLE);
\texttt{flash\_arc\_width} and \texttt{flash\_eff\_width} stable across
the swept range. Applicable to the blade data but not yet run there.

\textbf{Reynolds comparison} (\texttt{reynolds\_comparison.py}) ---
\textbf{CO₂ sweep only.} Defined with L = nozzle diameter d, U = bulk
gas velocity \texttt{v\_bulk\ =\ Q\ /\ (π/4\ ·\ d²)} (\emph{not} arc
velocity), \texttt{Re\ =\ ρ\ v\_bulk\ d\ /\ μ}, CO₂ at 295.15 K /
101.325 kPa (ρ = 1.815 kg·m⁻³, CRC 90th; μ = 1.48 × 10⁻⁵ Pa·s, Lemmon \&
Jacobsen 2004). Repeat-aware error bars on both axes. Caveats stated in
every caption: (i) changing d changes Re \textbf{and} geometry
(confound); (ii) Ma reaches ≈ 0.74 at the small-nozzle corner, so the
reported Re is incompressible-equivalent. Coverage (median arc
velocity): Re ≈ 3.3×10³ → 4.9×10³ → 6.5×10³ → 9.8×10³ giving 4.21 → 7.24
→ 9.47 → 10.01 m/s. The N₂ blade is a \textbf{single} condition (Re ≈
3×10³, Ma ≈ 0.25), so no Reynolds \emph{trend} exists for it --- only
the single value, for context.

\textbf{Repeat handling} (\texttt{sweep\_utils.py}). The 0.4 mm/1500 mL
CO₂ condition has two recordings disagreeing by ≈ 3.9 m/s; a
matched-settings diagnostic showed +3.35 m/s of genuine run-to-run
physics (not threshold/calibration), so the repeat is excluded from
cross-recording aggregates via an \texttt{.excluded\_from\_sweep}
sentinel while its per-recording outputs remain valid.

\textbf{Temporal drift \& long recordings.}
\texttt{temporal\_velocity\_profiles.py} and \texttt{warmup\_time.py}
test within-recording drift. Long \texttt{.cine} files must be probed
(\texttt{long\_video\_probe.py}) before trusting a prefix --- the J 1.2
mm recording showed a stable +0.34 m/s·s⁻¹ drift only after full
reprocessing (669 events vs 74 in the prefix). The blade recordings are
long (291k / 459k frames) and only representative segments have been
analysed --- the same probe applies before any prefix is trusted.

\begin{center}\rule{0.5\linewidth}{0.5pt}\end{center}

\hypertarget{12-provenance--reproducibility-common}{%
\subsection{12. Provenance \& reproducibility
(common)}\label{12-provenance--reproducibility-common}}

\texttt{run\_metadata.json} records, per run: fps, px/mm + source,
\texttt{bright\_thresh} + source, geometry mode + finite-segment flags,
gap\_row, \texttt{local\_vel\_n}, the σ\_pxmm/σ\_fps budget, git hash,
and UTC timestamp. \texttt{test\_j\_geometry.py} carries the
geometry-derivation and fps dimensional self-checks
(pytest-discoverable). \texttt{verify\_fps\_threading.py} audits that
the probed fps threads downstream across recordings.

\begin{center}\rule{0.5\linewidth}{0.5pt}\end{center}

\hypertarget{13-parameter-reference-side-by-side}{%
\subsection{13. Parameter reference (side by
side)}\label{13-parameter-reference-side-by-side}}

\begin{longtable}[]{@{}lll@{}}
\toprule
Parameter & CO₂ J-electrode & N₂ blade\tabularnewline
\midrule
\endhead
Detection threshold & per-recording (CLI \textgreater{} json
\textgreater{} txt \textgreater{} \textbf{210}) & \textbf{100 DN}
(dataset constant)\tabularnewline
px/mm & per-recording (≈ 2.4 mm apex gap); fallback 9.167 & \textbf{8.8}
(44 px = 5 mm)\tabularnewline
Electrode model & finite J-segments + \texttt{-\/-finite-segments} gate
& infinite blade lines + perp-distance test\tabularnewline
Geometry file & \texttt{j\_geometry.json} +
\texttt{blade\_lines.npy}/\texttt{electrode\_edges.npz} &
\texttt{blade\_lines.npy}\tabularnewline
fps & per-recording, auto-detected & 200 000 (5 µs/frame)\tabularnewline
\texttt{local\_vel\_n} & 10 (stable ≥ 7) & 10\tabularnewline
Mode A/B threshold & 0.15 & 0.15\tabularnewline
σ\_pxmm / σ\_fps & 5 \% / 0.1 \% (override per-recording) & 5 \% / 0.1
\%\tabularnewline
Operating points & sweep (4 conditions + 1 repeat) & single (2
orientations)\tabularnewline
\bottomrule
\end{longtable}

\begin{center}\rule{0.5\linewidth}{0.5pt}\end{center}

\hypertarget{14-known-limitations-common-unless-noted}{%
\subsection{14. Known limitations (common unless
noted)}\label{14-known-limitations-common-unless-noted}}

\begin{itemize}
\tightlist
\item
  \textbf{Motion blur} at the few-µs exposure inflates FWHM width ∝ tip
  speed.
\item
  \textbf{Single-distance calibration}; no lens-distortion correction
  (the 5 \% px/mm floor reflects this).
\item
  \textbf{Per-frame normalisation} (display/legacy) means cross-frame
  brightness must use raw-DN \texttt{integrated\_intensity}.
\item
  \textbf{CO₂-specific:} Mach \textgreater{} 0.3 at the small-nozzle
  corner (incompressible- equivalent Re); 0.4 mm/1500 mL effectively
  single-recording after the repeat exclusion; the 1.2 mm
  \texttt{j\_geometry.json} px/mm field is stale (8.75 vs 11.17).
\item
  \textbf{Blade-specific:} single operating condition (no Reynolds
  trend); only representative segments analysed so far; "OA" in the
  filename not yet expanded.
\item
  \textbf{Both:} Tier-B (re-detection) threshold sweep not run.
\end{itemize}

\begin{center}\rule{0.5\linewidth}{0.5pt}\end{center}

\hypertarget{15-references}{%
\subsection{15. References}\label{15-references}}

\begin{enumerate}
\def\labelenumi{\arabic{enumi}.}
\tightlist
\item
  Czernichowski, A. (1994). \emph{Pure Appl. Chem.} \textbf{66}(6),
  1301--1310.
\item
  Zhang, H., et al. (2020). \emph{IET High Voltage} \textbf{5}(4),
  386--392.
\item
  Fridman, A., Gutsol, A., \& Kennedy, L. A. (2005).
  \emph{Non-Equilibrium Air Plasmas at Atmospheric Pressure.} IOP
  Publishing.
\item
  Otsu, N. (1979). \emph{IEEE Trans. Syst. Man Cybern.} \textbf{9}(1),
  62--66.
\item
  Zhang, T. Y., \& Suen, C. Y. (1984). \emph{Commun. ACM}
  \textbf{27}(3), 236--239.
\item
  van der Walt, S. et al. (2014). "scikit-image." \emph{PeerJ}
  \textbf{2}, e453.
\item
  Lemmon, E. W., \& Jacobsen, R. T. (2004). \emph{Int. J. Thermophys.}
  \textbf{25}, 21--69 (gas transport properties).
\item
  CRC Handbook of Chemistry and Physics, 90th ed. (2009). CRC Press (gas
  densities).
\end{enumerate}

\begin{center}\rule{0.5\linewidth}{0.5pt}\end{center}

\emph{End of document.}

\end{document}
